\documentclass[aspectratio=169,11pt]{beamer}

\usetheme{AnnArbor}
\usecolortheme{dolphin}
\setbeamertemplate{navigation symbols}{}
\setbeamertemplate{footline}[frame number]

\usepackage{lmodern}
\usepackage[T1]{fontenc}

\usepackage{booktabs}
\usepackage{array}
\usepackage{tabularx}
\usepackage{listings}
\usepackage{xcolor}
\usepackage{amsmath}

\definecolor{codegray}{RGB}{245,245,245}
\definecolor{codeblue}{RGB}{0,65,120}
\definecolor{darkgreen}{RGB}{0,100,60}

\lstdefinestyle{pythonstyle}{
  language=Python,
  basicstyle=\ttfamily\scriptsize,
  keywordstyle=\color{codeblue}\bfseries,
  stringstyle=\color{darkgreen},
  commentstyle=\color{gray},
  showstringspaces=false,
  breaklines=true,
  frame=single,
  backgroundcolor=\color{codegray},
  columns=fullflexible
}

\lstdefinestyle{textstyle}{
  basicstyle=\ttfamily\scriptsize,
  showstringspaces=false,
  breaklines=true,
  frame=single,
  backgroundcolor=\color{codegray},
  columns=fullflexible
}

\lstset{style=pythonstyle}

\newcommand{\code}[1]{\texttt{#1}}
\newcommand{\smallgap}{\vspace{0.4em}}

\title[Data Access and Storage]{Data Access and Storage with Python}
\subtitle{Accessing, Importing, Integrating, and Storing Data for Analysis}
\author{DSAI1005 - Data Analysis with Python}
\institute{Faculty of Data Science and Artificial Intelligence\\National Economics University}
\date{2026}

\AtBeginSection[]{
  \begin{frame}{Roadmap}
    \tableofcontents[currentsection]
  \end{frame}
}

\begin{document}

\begin{frame}
  \titlepage
\end{frame}

\begin{frame}{Lesson Motivation}
In Data Science and Machine Learning projects, data is rarely available as a clean and analysis-ready DataFrame.

\smallgap
Data may be stored in many formats:
\begin{itemize}
  \item CSV and delimited text files
  \item Excel workbooks
  \item JSON returned by Web APIs
  \item HTML tables on websites
  \item PDF reports and documents
  \item relational databases and NoSQL databases
\end{itemize}
\end{frame}

\begin{frame}{Where This Topic Fits in the ML Workflow}
\begin{columns}[T]
\begin{column}{0.45\textwidth}
\begin{center}
\small
\begin{tabular}{c}
Define Purpose \\
$\downarrow$ \\
\textbf{Obtain Data} \\
$\downarrow$ \\
\textbf{Explore and Clean Data} \\
$\downarrow$ \\
Determine ML Task \\
$\downarrow$ \\
Choose ML Methods \\
$\downarrow$ \\
Train and Evaluate \\
$\downarrow$ \\
Deploy
\end{tabular}
\end{center}
\end{column}
\begin{column}{0.50\textwidth}
This lesson focuses on the work that must happen before modeling begins:
\begin{itemize}
  \item obtaining data;
  \item preparing data;
  \item integrating sources;
  \item storing reusable results.
\end{itemize}
\end{column}
\end{columns}
\end{frame}

\begin{frame}{Learning Objectives}
After this lesson, learners should be able to:
\begin{itemize}
  \item access data from CSV, Excel, JSON, HTML, PDF, SQLite, and MongoDB;
  \item use key Pandas I/O parameters such as \code{usecols}, \code{dtype}, \code{parse\_dates}, and \code{chunksize};
  \item flatten nested JSON data into tabular form;
  \item move data between SQL databases and Pandas DataFrames;
  \item build a multi-source data pipeline;
  \item choose suitable storage formats for analytical tasks.
\end{itemize}
\end{frame}

\begin{frame}{Lesson Structure}
\begin{enumerate}
  \item Overview of data access and storage
  \item CSV files with NumPy
  \item CSV files with Pandas
  \item Microsoft Excel
  \item JSON
  \item HTML and PDF
  \item SQLite
  \item MongoDB and PyMongo
  \item Multi-source data pipelines
  \item Best practices
\end{enumerate}
\end{frame}

\begin{frame}{Prerequisites}
To follow the lesson effectively, learners should have:
\begin{itemize}
  \item basic Python knowledge;
  \item familiarity with variables, lists, dictionaries, loops, and functions;
  \item basic knowledge of NumPy and Pandas DataFrames;
  \item understanding of rows, columns, and tabular data;
  \item access to Jupyter Notebook, Google Colab, VS Code, or a similar environment.
\end{itemize}
\end{frame}

\section{Overview of Data Access and Storage}

\begin{frame}[fragile]{What Is Data Access?}
\textbf{Data access} is the process of retrieving data from a storage source and bringing it into an analytical environment.

\smallgap
Example: loading a CSV file into Pandas.
\begin{lstlisting}[style=textstyle]
customers.csv
      |
      v
pd.read_csv()
      |
      v
DataFrame
\end{lstlisting}

Example: querying a database before analysis.
\begin{lstlisting}[style=textstyle]
SQL Database -> SQL Query -> Pandas DataFrame
\end{lstlisting}
\end{frame}

\begin{frame}[fragile]{What Is Data Storage?}
\textbf{Data storage} is the process of saving data or analytical results in a reusable format or storage system.

\begin{lstlisting}[style=textstyle]
DataFrame -> to_csv() -> report.csv
DataFrame -> to_sql() -> Database Table
\end{lstlisting}

\smallgap
Storage is not only an output step. It also supports reproducibility, reporting, sharing, and later modeling.
\end{frame}

\begin{frame}{Common Data Sources}
\begin{tabularx}{\textwidth}{lX}
\toprule
\textbf{Source Type} & \textbf{Examples} \\
\midrule
Flat files & CSV, TSV \\
Spreadsheets & Excel \\
Semi-structured data & JSON \\
Web data & HTML tables \\
Documents & PDF reports \\
Relational databases & SQLite, PostgreSQL, MySQL \\
NoSQL databases & MongoDB \\
\bottomrule
\end{tabularx}
\end{frame}

\begin{frame}[fragile]{Structured and Semi-Structured Data}
Structured data has a clear row-and-column structure.
\begin{lstlisting}[style=textstyle]
CustomerID | Name | City | Revenue
\end{lstlisting}

Semi-structured data may contain nested objects and lists.
\begin{lstlisting}[style=pythonstyle]
{
    "customer": {"name": "An", "city": "Hanoi"},
    "orders": [
        {"product": "A"},
        {"product": "B"}
    ]
}
\end{lstlisting}
\end{frame}

\begin{frame}{Exercise: Classify Data Sources}
Classify each source into a suitable category.
\begin{enumerate}
  \item \code{sales.csv}
  \item \code{customers.xlsx}
  \item JSON data returned by a Web API
  \item \code{company.db}
  \item MongoDB
  \item A financial report in PDF format
\end{enumerate}

\smallgap
Use: flat file, spreadsheet, semi-structured data, relational database, NoSQL database, document.
\end{frame}

\begin{frame}{Knowledge Check}
\textbf{Question 1.} Which format is commonly considered a flat file?
\begin{enumerate}[A.]
  \item CSV
  \item MongoDB
  \item SQLite
  \item PDF
\end{enumerate}

\smallgap
\textbf{Question 2.} Which format naturally supports nested objects and lists?
\begin{enumerate}[A.]
  \item JSON
  \item CSV
  \item TSV
  \item Numerical matrix
\end{enumerate}
\end{frame}

\section{Working with CSV Files Using NumPy}

\begin{frame}[fragile]{Reading Numerical Files with \code{np.loadtxt()}}
\code{np.loadtxt()} is suitable when data is mainly numerical, simple, and has no missing values.

\begin{lstlisting}[style=pythonstyle]
import numpy as np

data = np.loadtxt(
    "matrix_data.csv",
    delimiter=",",
    skiprows=1
)

print(data)
print(data.shape)
\end{lstlisting}

Column means can be computed using:
\begin{lstlisting}[style=pythonstyle]
print(np.mean(data, axis=0))
\end{lstlisting}
\end{frame}

\begin{frame}[fragile]{Exercise: \code{np.loadtxt()}}
Complete the program to read \code{sales.csv}, where values are separated by commas.

\begin{lstlisting}[style=pythonstyle]
# data = np.loadtxt(
#     "sales.csv",
#     delimiter=...
# )

# print(data.shape)
\end{lstlisting}

Then determine the number of rows and columns.
\end{frame}

\begin{frame}[fragile]{Handling Missing Values with \code{np.genfromtxt()}}
\code{np.genfromtxt()} is more flexible than \code{np.loadtxt()} when numerical files contain missing values.

\begin{lstlisting}[style=pythonstyle]
data = np.genfromtxt(
    "numeric_data.csv",
    delimiter=",",
    skip_header=1,
    filling_values=0.0,
    dtype=float
)

print(data)
\end{lstlisting}

Here, missing values are replaced with \code{0.0}.
\end{frame}

\begin{frame}[fragile]{Writing Arrays with \code{np.savetxt()}}
\code{np.savetxt()} saves a NumPy array to a text file.

\begin{lstlisting}[style=pythonstyle]
output = np.random.randn(5, 3)

np.savetxt(
    "output.csv",
    output,
    delimiter=",",
    fmt="%.2f",
    header="A,B,C",
    comments=""
)
\end{lstlisting}

\code{fmt="\%.2f"} writes each number with two decimal places.
\end{frame}

\begin{frame}{NumPy CSV Functions}
\begin{tabularx}{\textwidth}{lX}
\toprule
\textbf{Function} & \textbf{Purpose} \\
\midrule
\code{np.loadtxt()} & Read numerical data from simple structured text files \\
\code{np.genfromtxt()} & Read numerical data with missing-value support \\
\code{np.savetxt()} & Write NumPy arrays to text files \\
\bottomrule
\end{tabularx}
\end{frame}

\begin{frame}[fragile]{Practical Exercise: Monthly Sales}
Create \code{monthly\_sales.csv} with the following content:
\begin{lstlisting}[style=textstyle]
Month,Revenue,Cost
1,120,80
2,150,
3,180,110
4,,130
\end{lstlisting}

Complete these tasks:
\begin{enumerate}
  \item Read the data using \code{np.genfromtxt()}.
  \item Replace missing values with \code{0}.
  \item Calculate mean Revenue and Cost.
  \item Save the processed numerical data using \code{np.savetxt()}.
\end{enumerate}
\end{frame}

\section{Working with CSV Files Using Pandas}

\begin{frame}[fragile]{Reading CSV Files with \code{pd.read\_csv()}}
Pandas provides \code{pd.read\_csv()} for reading tabular data from CSV files.

\begin{lstlisting}[style=pythonstyle]
import pandas as pd

df = pd.read_csv("customers.csv")

print(df.head())
\end{lstlisting}

Pandas is suitable when data contains text, dates, multiple types, missing values, and column labels.
\end{frame}

\begin{frame}[fragile]{Specifying Delimiters with \code{sep}}
Not every file uses a comma as the separator.

Semicolon-delimited file:
\begin{lstlisting}[style=pythonstyle]
df = pd.read_csv(
    "data.csv",
    sep=";"
)
\end{lstlisting}

Tab-separated file:
\begin{lstlisting}[style=pythonstyle]
df = pd.read_csv(
    "data.tsv",
    sep="\t"
)
\end{lstlisting}
\end{frame}

\begin{frame}[fragile]{Reading Selected Columns with \code{usecols}}
\code{usecols} selects columns during import.

\begin{lstlisting}[style=pythonstyle]
df = pd.read_csv(
    "customers.csv",
    usecols=[
        "customer_id",
        "city",
        "total_spent"
    ]
)
\end{lstlisting}

Benefits include lower memory use, faster loading, and simpler downstream analysis.
\end{frame}

\begin{frame}[fragile]{Specifying Data Types with \code{dtype}}
Pandas can infer data types automatically, but explicit types are often safer.

\begin{lstlisting}[style=pythonstyle]
df = pd.read_csv(
    "customers.csv",
    dtype={
        "customer_id": str,
        "age": "Int64"
    }
)
\end{lstlisting}

Fields such as CustomerID, StudentID, ProductID, and PostalCode usually represent identifiers, not numerical measurements.
\end{frame}

\begin{frame}[fragile]{Parsing Dates with \code{parse\_dates}}
Pandas can convert columns to datetime during import.

\begin{lstlisting}[style=pythonstyle]
df = pd.read_csv(
    "orders.csv",
    parse_dates=["OrderDate"]
)

print(df.dtypes)
\end{lstlisting}

This is useful before filtering, grouping, or resampling by time.
\end{frame}

\begin{frame}[fragile]{Defining Missing Values with \code{na\_values}}
Real datasets may use different symbols for missing values.

\begin{lstlisting}[style=pythonstyle]
df = pd.read_csv(
    "customers.csv",
    na_values=[
        "NA",
        "N/A",
        "-",
        "missing"
    ]
)
\end{lstlisting}

Pandas converts these symbols into standard missing values.
\end{frame}

\begin{frame}[fragile]{Character Encoding}
Use UTF-8 when reading text with non-ASCII characters.

\begin{lstlisting}[style=pythonstyle]
df = pd.read_csv(
    "customers.csv",
    encoding="utf-8"
)
\end{lstlisting}

For CSV files intended for Microsoft Excel on Windows:
\begin{lstlisting}[style=pythonstyle]
df.to_csv(
    "customers_clean.csv",
    index=False,
    encoding="utf-8-sig"
)
\end{lstlisting}
\end{frame}

\begin{frame}[fragile]{Inspecting Data After Import}
After loading a dataset, inspect structure and quality immediately.

\begin{lstlisting}[style=pythonstyle]
df.head()
df.shape
df.info()
df.dtypes
df.isna().sum()
\end{lstlisting}

These commands help identify row counts, column names, data types, non-missing values, and missing-value counts.
\end{frame}

\begin{frame}[fragile]{Reading Large Datasets with \code{chunksize}}
When a CSV file is larger than available memory, process it incrementally.

\begin{lstlisting}[style=pythonstyle]
total = 0

for chunk in pd.read_csv(
    "transactions.csv",
    chunksize=100000
):
    total += chunk["Amount"].sum()

print(total)
\end{lstlisting}

The file is read, processed, and discarded one chunk at a time.
\end{frame}

\begin{frame}[fragile]{Writing a DataFrame to CSV}
Use \code{to\_csv()} to export a DataFrame.

\begin{lstlisting}[style=pythonstyle]
df.to_csv(
    "clean_data.csv",
    index=False,
    encoding="utf-8-sig"
)
\end{lstlisting}

\code{index=False} prevents the DataFrame index from being written as an additional column.
\end{frame}

\begin{frame}{Pandas CSV Functions and Arguments}
\begin{tabularx}{\textwidth}{lX}
\toprule
\textbf{Function or Argument} & \textbf{Purpose} \\
\midrule
\code{pd.read\_csv()} & Read CSV or other delimited files \\
\code{usecols=} & Read selected columns only \\
\code{dtype=} & Specify data types \\
\code{parse\_dates=} & Convert columns to datetime \\
\code{na\_values=} & Define missing-value symbols \\
\code{chunksize=} & Read data incrementally \\
\code{df.to\_csv()} & Write a DataFrame to CSV \\
\bottomrule
\end{tabularx}
\end{frame}

\begin{frame}{Knowledge Check: Pandas CSV}
\textbf{Question 1.} Which argument allows only selected columns to be read?
\begin{enumerate}[A.]
  \item \code{usecols}
  \item \code{keepcols}
  \item \code{columns\_only}
  \item \code{filter\_columns}
\end{enumerate}

\smallgap
\textbf{Question 2.} Which argument is useful for very large CSV files?
\begin{enumerate}[A.]
  \item \code{chunksize}
  \item \code{batch=True}
  \item \code{split=True}
  \item \code{large=True}
\end{enumerate}
\end{frame}

\begin{frame}{Practical Exercise: Customer CSV}
Create \code{customers.csv} containing CustomerID, Name, City, Age, SignupDate, and TotalSpent.

Include missing symbols such as \code{NA} and \code{-}, and CustomerID values such as \code{001} and \code{002}.

Tasks:
\begin{enumerate}
  \item Read CustomerID as \code{str}.
  \item Parse SignupDate as datetime.
  \item Treat \code{NA} and \code{-} as missing values.
  \item Read only five columns needed for analysis.
  \item Inspect the data and export \code{customers\_clean.csv}.
\end{enumerate}
\end{frame}

\section{Working with Microsoft Excel}

\begin{frame}[fragile]{Reading a Worksheet}
Pandas uses \code{pd.read\_excel()} to read Excel data.

\begin{lstlisting}[style=pythonstyle]
df = pd.read_excel(
    "products.xlsx",
    sheet_name="Products",
    engine="openpyxl"
)
\end{lstlisting}

The \code{sheet\_name} argument specifies the worksheet to read.
\end{frame}

\begin{frame}[fragile]{Reading All Worksheets}
An entire workbook can be read by setting \code{sheet\_name=None}.

\begin{lstlisting}[style=pythonstyle]
workbook = pd.read_excel(
    "products.xlsx",
    sheet_name=None
)

print(workbook.keys())
products = workbook["Products"]
\end{lstlisting}

The result is a dictionary of DataFrames.
\end{frame}

\begin{frame}[fragile]{Writing Multiple DataFrames to One Workbook}
Use \code{pd.ExcelWriter()} to write multiple DataFrames to separate worksheets.

\begin{lstlisting}[style=pythonstyle]
with pd.ExcelWriter(
    "report.xlsx",
    engine="openpyxl"
) as writer:
    df_products.to_excel(
        writer,
        sheet_name="Products",
        index=False
    )
    df_sales.to_excel(
        writer,
        sheet_name="Sales",
        index=False
    )
\end{lstlisting}
\end{frame}

\begin{frame}{Excel Functions}
\begin{tabularx}{\textwidth}{lX}
\toprule
\textbf{Function or Argument} & \textbf{Purpose} \\
\midrule
\code{pd.read\_excel()} & Read Excel data \\
\code{sheet\_name=} & Select a worksheet \\
\code{sheet\_name=None} & Read all worksheets \\
\code{df.to\_excel()} & Write a DataFrame to Excel \\
\code{pd.ExcelWriter()} & Write multiple DataFrames to one workbook \\
\bottomrule
\end{tabularx}
\end{frame}

\begin{frame}{Excel Exercise and Knowledge Check}
\textbf{Exercise.} Create \code{business\_data.xlsx} with three worksheets: Customers, Products, and Sales. Then read the workbook, inspect worksheet names and dimensions, calculate total Sales, and export a summary workbook.

\smallgap
\textbf{Question 1.} What does \code{pd.read\_excel()} return when \code{sheet\_name=None} is used?
\begin{enumerate}[A.]
  \item A dictionary of DataFrames
  \item Only the first worksheet
  \item A NumPy array
  \item A string
\end{enumerate}
\end{frame}

\section{Working with JSON}

\begin{frame}[fragile]{What Is JSON?}
JSON, or JavaScript Object Notation, is widely used in Web APIs, e-commerce systems, microservices, and NoSQL databases.

\begin{lstlisting}[style=textstyle]
[
    {"customer_id": "C001", "city": "Hanoi"},
    {"customer_id": "C002", "city": "Danang"}
]
\end{lstlisting}

For relatively flat JSON structures:
\begin{lstlisting}[style=pythonstyle]
df = pd.read_json("customers.json")
\end{lstlisting}
\end{frame}

\begin{frame}[fragile]{Nested JSON}
Real-world JSON often contains multiple levels.

\begin{lstlisting}[style=pythonstyle]
{
    "order_id": "O001",
    "customer": {
        "customer_id": "C001",
        "city": "Hanoi"
    },
    "items": [
        {"sku": "P001", "qty": 2},
        {"sku": "P002", "qty": 1}
    ]
}
\end{lstlisting}

For analysis, each item may need to become a separate row.
\end{frame}

\begin{frame}[fragile]{Flattening Data with \code{pd.json\_normalize()}}
\begin{lstlisting}[style=pythonstyle]
df_items = pd.json_normalize(
    orders,
    record_path=["items"],
    meta=[
        "order_id",
        ["customer", "customer_id"],
        ["customer", "city"]
    ]
)
\end{lstlisting}

\begin{itemize}
  \item \code{record\_path} specifies the nested list to expand into rows.
  \item \code{meta} specifies parent-level attributes to retain.
\end{itemize}
\end{frame}

\begin{frame}{JSON Exercise and Knowledge Check}
\textbf{Exercise.} Create three orders in JSON. Each order contains OrderID, OrderDate, Customer, and Items. Flatten the data so that each item becomes one row, then calculate Revenue by City.

\smallgap
\textbf{Question.} What is the purpose of \code{record\_path} in \code{pd.json\_normalize()}?
\begin{enumerate}[A.]
  \item To identify the nested list to expand into rows
  \item To rename columns
  \item To write data to JSON
  \item To remove missing values
\end{enumerate}
\end{frame}

\section{Accessing Data from HTML and PDF}

\begin{frame}[fragile]{Reading HTML Tables}
\code{pd.read\_html()} detects HTML \code{<table>} elements and converts them into DataFrames.

\begin{lstlisting}[style=pythonstyle]
tables = pd.read_html(
    "financial_quotes.html"
)

print(len(tables))
stocks = tables[0]
\end{lstlisting}

The result is a list of DataFrames.
\end{frame}

\begin{frame}{Why PDF Extraction Is Harder}
CSV is designed to represent tabular data.

PDF is primarily designed to:
\begin{itemize}
  \item display content;
  \item preserve layout;
  \item support printing.
\end{itemize}

Rows and columns that appear visually may not be stored internally as explicit table structures.
\end{frame}

\begin{frame}[fragile]{Extracting Tables with \code{pdfplumber}}
\begin{lstlisting}[style=pythonstyle]
import pdfplumber

with pdfplumber.open("report.pdf") as pdf:
    page = pdf.pages[0]
    tables = page.extract_tables()
\end{lstlisting}

Convert one extracted table to a DataFrame:
\begin{lstlisting}[style=pythonstyle]
table = tables[0]

df = pd.DataFrame(
    table[1:],
    columns=table[0]
)
\end{lstlisting}
\end{frame}

\begin{frame}{HTML and PDF Exercise}
Given an HTML file containing a product-price table and a foreign-exchange table:
\begin{enumerate}
  \item read both tables into Pandas;
  \item inspect the structure of each table;
  \item save each table to a separate worksheet in one Excel workbook.
\end{enumerate}

\smallgap
For a table extracted from a PDF, inspect headers, missing values, and required cleaning operations before analysis.
\end{frame}

\section{Working with SQLite Databases}

\begin{frame}[fragile]{What Is SQLite?}
SQLite is a lightweight relational database system. An entire database can be stored in one file, such as \code{company.db}.

Python provides SQLite support through the standard library.
\begin{lstlisting}[style=pythonstyle]
import sqlite3
\end{lstlisting}

Create a connection and cursor:
\begin{lstlisting}[style=pythonstyle]
conn = sqlite3.connect("company.db")
cursor = conn.cursor()
\end{lstlisting}
\end{frame}

\begin{frame}[fragile]{Executing SQL Queries}
The following query retrieves employees with salary at least 30 million.

\begin{lstlisting}[style=pythonstyle]
cursor.execute(
    """
    SELECT emp_id, full_name, salary
    FROM employees
    WHERE salary >= ?
    """,
    (30000000,)
)

rows = cursor.fetchall()
\end{lstlisting}

\code{fetchall()} returns all rows satisfying the query.
\end{frame}

\begin{frame}[fragile]{Parameterized Queries}
When query values come from variables or user input, use parameters.

\begin{lstlisting}[style=pythonstyle]
cursor.execute(
    "SELECT * FROM users WHERE id = ?",
    (user_id,)
)
\end{lstlisting}

This separates SQL structure from values and helps reduce security risks associated with direct string concatenation.
\end{frame}

\begin{frame}[fragile]{Reading SQL Results into Pandas}
Pandas can read SQL query results directly into a DataFrame.

\begin{lstlisting}[style=pythonstyle]
query = """
SELECT
    emp_id,
    full_name,
    salary
FROM employees
"""

df = pd.read_sql_query(
    query,
    conn
)
\end{lstlisting}
\end{frame}

\begin{frame}[fragile]{Combining Tables with JOIN}
\begin{lstlisting}[style=textstyle]
SELECT
    e.emp_id,
    e.full_name,
    d.dept_name
FROM employees e
LEFT JOIN departments d
ON e.dept_id = d.dept_id;
\end{lstlisting}

The query result can then be loaded into Pandas.
\begin{lstlisting}[style=pythonstyle]
df = pd.read_sql_query(
    query,
    conn
)
\end{lstlisting}
\end{frame}

\begin{frame}[fragile]{Writing DataFrames to SQLite}
\begin{lstlisting}[style=pythonstyle]
df.to_sql(
    name="sales_summary",
    con=conn,
    if_exists="replace",
    index=False
)
\end{lstlisting}

Common values of \code{if\_exists}:
\begin{itemize}
  \item \code{fail}: raise an error if the table exists;
  \item \code{replace}: replace the existing table;
  \item \code{append}: add rows to the existing table.
\end{itemize}
\end{frame}

\begin{frame}{SQLite Functions and Exercise}
\begin{tabularx}{\textwidth}{lX}
\toprule
\textbf{Function} & \textbf{Purpose} \\
\midrule
\code{sqlite3.connect()} & Open a connection to SQLite \\
\code{cursor.execute()} & Execute an SQL statement \\
\code{fetchall()} & Retrieve query results \\
\code{pd.read\_sql\_query()} & Convert SQL results to DataFrame \\
\code{df.to\_sql()} & Write DataFrame to SQL table \\
\bottomrule
\end{tabularx}

\smallgap
\textbf{Exercise.} Create \code{sales.db}, join \code{orders} and \code{customers}, calculate total Amount by City, and store the result in \code{city\_sales\_summary}.
\end{frame}

\section{Working with MongoDB and PyMongo}

\begin{frame}[fragile]{What Is MongoDB?}
MongoDB is a document-oriented NoSQL database system.

\begin{columns}[T]
\begin{column}{0.45\textwidth}
Relational database:
\begin{lstlisting}[style=textstyle]
Table
  -> Row
\end{lstlisting}
\end{column}
\begin{column}{0.45\textwidth}
MongoDB:
\begin{lstlisting}[style=textstyle]
Collection
  -> Document
\end{lstlisting}
\end{column}
\end{columns}

\smallgap
Documents can contain scalar values, nested objects, lists, and multiple levels of hierarchy.
\end{frame}

\begin{frame}[fragile]{Connecting to MongoDB}
PyMongo provides \code{MongoClient} for connecting to MongoDB.

\begin{lstlisting}[style=pythonstyle]
from pymongo import MongoClient

client = MongoClient(
    "mongodb://localhost:27017/"
)

db = client["retail_db"]
orders = db["orders"]
\end{lstlisting}
\end{frame}

\begin{frame}[fragile]{CRUD Operations}
CRUD refers to Create, Read, Update, and Delete.

\begin{lstlisting}[style=pythonstyle]
# Create
orders.insert_one(sample_order)

# Read one document
orders.find_one({"order_id": "ORD001"})

# Filter multiple documents
orders.find({"total_amount": {"$gte": 1000000}})

# Update
orders.update_one(
    {"order_id": "ORD001"},
    {"$set": {"status": "REFUNDED"}}
)
\end{lstlisting}
\end{frame}

\begin{frame}[fragile]{Converting MongoDB Data to Pandas}
Query results can be converted into a list of documents.

\begin{lstlisting}[style=pythonstyle]
documents = list(
    orders.find()
)

df = pd.json_normalize(
    documents
)
\end{lstlisting}

MongoDB's \code{\_id} field can be converted before exporting.
\begin{lstlisting}[style=pythonstyle]
df["_id"] = df["_id"].astype(str)
\end{lstlisting}
\end{frame}

\begin{frame}{MongoDB Exercise and Knowledge Check}
\textbf{Exercise.} Create a collection named \code{customer\_orders} with at least five documents. Insert orders, query paid orders, update one status, convert results to Pandas, and calculate paid revenue.

\smallgap
\textbf{Question.} Which Pandas function is useful for nested MongoDB documents?
\begin{enumerate}[A.]
  \item \code{pd.json\_normalize()}
  \item \code{pd.read\_csv()}
  \item \code{pd.crosstab()}
  \item \code{pd.cut()}
\end{enumerate}
\end{frame}

\section{Building a Multi-Source Data Pipeline}

\begin{frame}[fragile]{Pipeline Scenario}
A company stores data in multiple sources:
\begin{lstlisting}[style=textstyle]
customers.csv
products.xlsx
orders.json
company.db
\end{lstlisting}

The goal is to integrate these sources into one analytical dataset for reporting, visualization, or Machine Learning.
\end{frame}

\begin{frame}{General Integration Workflow}
\begin{center}
\begin{tabular}{c}
CSV / Excel / JSON / SQL \\
$\downarrow$ \\
Data Ingestion \\
$\downarrow$ \\
Data Cleaning \\
$\downarrow$ \\
Transformation \\
$\downarrow$ \\
Merge \\
$\downarrow$ \\
Analytics Dataset \\
$\swarrow \quad \searrow$ \\
SQLite \qquad Excel Report
\end{tabular}
\end{center}
\end{frame}

\begin{frame}[fragile]{Reading Multiple Sources}
\begin{lstlisting}[style=pythonstyle]
customers = pd.read_csv("customers.csv")

products = pd.read_excel(
    "products.xlsx",
    sheet_name="Products"
)

import json
with open("orders.json", "r", encoding="utf-8") as f:
    orders = json.load(f)
\end{lstlisting}
\end{frame}

\begin{frame}[fragile]{Flattening Orders and Creating Variables}
\begin{lstlisting}[style=pythonstyle]
items = pd.json_normalize(
    orders,
    record_path=["items"],
    meta=[
        "order_id",
        ["customer", "customer_id"]
    ]
)

items["Revenue"] = (
    items["Quantity"] * items["Price"]
)
\end{lstlisting}
\end{frame}

\begin{frame}[fragile]{Integrating Data with \code{merge()}}
\begin{lstlisting}[style=pythonstyle]
df = items.merge(
    products,
    on="ProductID",
    how="left"
)

df = df.merge(
    customers,
    on="CustomerID",
    how="left"
)
\end{lstlisting}

After each merge, inspect:
\begin{lstlisting}[style=pythonstyle]
print(df.shape)
print(df.isna().sum())
\end{lstlisting}
\end{frame}

\begin{frame}[fragile]{Creating KPI Summaries}
Calculate revenue and number of orders by region.

\begin{lstlisting}[style=pythonstyle]
region_summary = (
    df
    .groupby("Region")
    .agg(
        Revenue=("Revenue", "sum"),
        Orders=("OrderID", "nunique")
    )
    .reset_index()
)
\end{lstlisting}
\end{frame}

\begin{frame}[fragile]{Storing and Reporting Results}
Store analytical data in SQLite.
\begin{lstlisting}[style=pythonstyle]
conn = sqlite3.connect("analytics.db")

df.to_sql(
    "order_analytics",
    conn,
    if_exists="replace",
    index=False
)

conn.close()
\end{lstlisting}

Export KPI summaries to Excel.
\begin{lstlisting}[style=pythonstyle]
with pd.ExcelWriter("dashboard.xlsx") as writer:
    df.to_excel(writer, sheet_name="Details", index=False)
    region_summary.to_excel(writer, sheet_name="Region_KPI", index=False)
\end{lstlisting}
\end{frame}

\begin{frame}{Project Exercise: Multi-Source Integration}
Build \code{order\_analytics} from \code{customers.csv}, \code{products.xlsx}, and \code{orders.json}.

Required fields:
\begin{itemize}
  \item OrderID, CustomerID, CustomerName, City
  \item ProductID, Category, Quantity, UnitPrice, Revenue
\end{itemize}

Tasks: read sources, inspect data, flatten JSON, merge tables, identify unmatched records, calculate revenue, save details to SQLite, and export KPI summaries to Excel.
\end{frame}

\section{Best Practices}

\begin{frame}[fragile]{Read Only Required Columns}
If a file contains many columns but only a subset is needed, select columns during import.

\begin{lstlisting}[style=pythonstyle]
pd.read_csv(
    "large.csv",
    usecols=[
        "CustomerID",
        "Sales"
    ]
)
\end{lstlisting}

This reduces loading time and memory consumption.
\end{frame}

\begin{frame}[fragile]{Inspect Early and Often}
Inspect immediately after import.
\begin{lstlisting}[style=pythonstyle]
df.head()
df.shape
df.info()
df.isna().sum()
\end{lstlisting}

Inspect again after integration.
\begin{lstlisting}[style=pythonstyle]
df.isna().sum()
\end{lstlisting}

New missing values may appear when keys do not match across sources.
\end{frame}

\begin{frame}[fragile]{Control Data Types}
Inspect data types using:
\begin{lstlisting}[style=pythonstyle]
df.dtypes
\end{lstlisting}

Do not assume that every field containing digits is numerical. For example, \code{CustomerID = "001"} is an identifier and should often be stored as a string.
\end{frame}

\begin{frame}[fragile]{Use Context Managers}
When working with files, use \code{with} statements.

\begin{lstlisting}[style=pythonstyle]
with open("data.json", encoding="utf-8") as f:
    data = json.load(f)
\end{lstlisting}

For Excel reports:
\begin{lstlisting}[style=pythonstyle]
with pd.ExcelWriter("report.xlsx") as writer:
    summary.to_excel(writer, index=False)
\end{lstlisting}

Resources are closed properly after the operation is complete.
\end{frame}

\begin{frame}{Choose an Appropriate Storage Format}
\begin{tabularx}{\textwidth}{lX}
\toprule
\textbf{Requirement} & \textbf{Suitable Format} \\
\midrule
Exchange simple tabular data & CSV \\
Create multi-sheet reports & Excel \\
Exchange data through Web APIs & JSON \\
Store relational data and queries & SQL \\
Store flexible nested documents & MongoDB \\
Store large tabular datasets & Parquet \\
Extract website tables & HTML \\
Distribute fixed-layout reports & PDF \\
\bottomrule
\end{tabularx}
\end{frame}

\begin{frame}{Final Knowledge Check}
\textbf{Question 1.} Which technique is appropriate when a CSV file is larger than RAM?
\begin{enumerate}[A.]
  \item Reading the file in chunks
  \item Converting all values to strings
  \item Creating a second copy of the file
  \item Using only \code{head()}
\end{enumerate}

\smallgap
\textbf{Question 2.} Which format is especially common for exchanging data through Web APIs?
\begin{enumerate}[A.]
  \item JSON
  \item PDF
  \item XLS
  \item SQLite
\end{enumerate}
\end{frame}

\begin{frame}{Summary}
A general workflow for working with data from multiple sources:
\begin{center}
\small
\begin{tabular}{ccccccccc}
Source & $\rightarrow$ & Read & $\rightarrow$ & Inspect & $\rightarrow$ & Clean & $\rightarrow$ & Transform \\[0.4em]
 & & & & & & & & $\downarrow$ \\[0.4em]
Store / Report & $\leftarrow$ & Analyze & $\leftarrow$ & Integrate & $\leftarrow$ & & &
\end{tabular}
\end{center}
\smallgap
The goal is to transform raw, distributed data into reliable analytical data.
\end{frame}

\begin{frame}{Key Tools}
\begin{tabularx}{\textwidth}{lX}
\toprule
\textbf{Data Source} & \textbf{Python Tools} \\
\midrule
CSV & \code{pd.read\_csv()}, \code{to\_csv()} \\
NumPy text files & \code{loadtxt()}, \code{genfromtxt()}, \code{savetxt()} \\
Excel & \code{read\_excel()}, \code{ExcelWriter()} \\
JSON & \code{read\_json()}, \code{json\_normalize()} \\
HTML & \code{read\_html()} \\
PDF & \code{pdfplumber} \\
SQLite & \code{sqlite3}, \code{read\_sql\_query()}, \code{to\_sql()} \\
MongoDB & \code{pymongo}, \code{json\_normalize()} \\
\bottomrule
\end{tabularx}
\end{frame}

\begin{frame}{Closing Message}
Data access, preparation, and storage are not secondary technical steps.

\smallgap
They determine whether later analysis, visualization, and Machine Learning models can be trusted.

\smallgap
\begin{center}
\Large Correct data pipeline $\Rightarrow$ reliable analytical results
\end{center}
\end{frame}

\end{document}
